\paragraph{Conjugate-coordinate class: tree-level quadratic form of the level-one amplitudes (conditional; not closure).}
This paragraph supplements the preceding ones, in which the tree-level action of the amplitudes \(b^{\nu;\mu\rho}\) was not computed. Keep the Poisson limit, the leaf \(\mathbb{CP}^{1,2}\), \(X^a=T^a+A^a\), the base functions \(y^\mu\) and the notation of the preceding paragraphs, in units \(R=1\), and write \(D^a=\{T^a,\cdot\}\), \(\mathcal J_{ab}=\{M_{ab},\cdot\}\), \(\Box=D_bD^b\). The quartic term of~\eqref{eq:Smatrix} with \(C_{\mu\nu\lambda}=0\), read with \([\cdot,\cdot]\to i\{\cdot,\cdot\}\) and \(\operatorname{Tr}\to\int\mathrm{vol}\), is \(S=(4g_{\mathrm{YM}}^2)^{-1}\int\mathrm{vol}\,\{X^a,X^b\}\{X_a,X_b\}\), with field equation \(\{X^b,\{X_b,X_a\}\}=0\).

\emph{Background.} The Lie--Poisson relations give \(\{T_a,\{T^a,T^b\}\}=-3T^b\), an identity of the polynomial Poisson algebra of \(\mathfrak{so}(4,2)\) that does not depend on the sign \(s\). Hence \(T^a\) is not stationary for the quartic term alone. It is stationary for \(S+(2g_{\mathrm{YM}}^2)^{-1}\mu^2\int\mathrm{vol}\,X_aX^a\), whose field equation is \(\{X^b,\{X_b,X_a\}\}=\mu^2X_a\), exactly when \(\mu^2=-3\). A constant \(C_{\mu\nu\lambda}\) cannot take the place of this term: it would add \(C_{abc}M^{bc}\) to the field equation, and these functions are linearly independent of \(T_a=M_{a4}\) on the orbit. The bilinear term is not part of~\eqref{eq:Smatrix}; it is an input of this class with a coefficient fixed by stationarity, and a mass term is reported for the covariant backgrounds of~\cite{SperlingSteinacker2019}.

\emph{Quadratic form.} With this term, \(S^{(2)}=(2g_{\mathrm{YM}}^2)^{-1}\int\mathrm{vol}\,A^aL(A)_a\) with \(L(A)_a=-\Box A_a+D_a(D_bA^b)-2s\,\mathcal J_{ab}A^b+\mu^2A_a\), an identity verified in the polynomial Poisson algebra for both signs \(s\). Gauge covariance of the field equation about a stationary point gives \(L(\{\Lambda,T\})=0\) for every function \(\Lambda\), and the same check fails for \(\mu^2\neq-3\). The quadratic form is therefore exactly invariant under \(\delta A^a=\{\Lambda,T^a\}\), not only at a leading symbol.

\emph{Leading symbol.} Since \(D^ay^\mu=s\,m\,\eta^{a\mu}\) with \(m=M^{45}\) a function of \(y\) alone, the only part of \(L\) with two derivatives along \(y\) is \(-\Box A_a+D_a(D_bA^b)\). The operator \(\mathcal J_{ab}\) acts on \(y^\mu\) as a Lorentz rotation with a coefficient linear in \(y\), and the fibre derivatives of \(A\) and the derivatives of \(m\) enter with coefficients of relative order \(1/(|q|\ell)\). For \(A^a=c^a(y)F\) with \(F\) a fixed fibre function and \(|q|\ell\gg1\), then, \(L(A)_a=-m^2(\partial^2c_a-\partial_a\partial_bc^b)F\) plus terms with at most one derivative of \(c\), so that
\[
S^{(2)}\simeq\frac{1}{4g_{\mathrm{YM}}^2}\int\mathrm{vol}\;m^2\,\mathcal F_{ab}\mathcal F^{ab},\qquad
\mathcal F_{ab}=\partial_aA_b-\partial_bA_a ,
\]
with \(\partial_a\) the derivative along \(y\) at fixed fibre argument. At this order each fibre function is an independent Maxwell field with the same coefficient, and there is no coupling between fibre harmonics. On the explicit realisation of the preceding paragraphs, evaluating \(L\) on \(A^a=c^a\mathrm e^{i\lambda q\cdot y}M^{12}\) reproduces \(\lambda^2m^2M^{12}\mathrm e^{i\lambda q\cdot y}(q^2c_a-q_aq\cdot c)\) with a relative remainder that falls as \(1/\lambda\).

\emph{Level-one sector.} Write \(M^{\mu\rho}=\sum_ie^iS_i^{\mu\rho}\) with \(S_i\) an orthogonal, equal-norm basis of \(W_y\) and \(e^i\) the three first spherical harmonics on the fibre; by the \(\mathrm{SU}(2)\) invariance of the fibre measure, \(\int\mathrm d\mu\,e^ie^j=\varkappa(y)\,\delta^{ij}\) with \(\varkappa>0\). For \(A^\nu=c^\nu_i(y)e^i\) one has \(b^{\nu;\mu\rho}=\varkappa\,c^\nu_iS_i^{\mu\rho}\), and the twelve amplitudes are three vector fields \(c_i\) with the action \((m^2\varkappa/4g_{\mathrm{YM}}^2)\sum_i\int\mathrm d^4y\,\mathcal F_{i,ab}\mathcal F_i^{ab}\). The kinetic matrix is \(\delta_{ij}(q^2\eta_{ab}-q_aq_b)\), definite in the index \(i\); the leading symbol has no mass term and no coupling to the other harmonics. The overall sign is that of the quartic term as displayed, common to all fibre harmonics including the fibre-constant gauge field, so no relative ghost arises at this order. It is opposite to the sign of the Lorentzian Maxwell action \(-\tfrac14F_{ab}F^{ab}\) in the same \(\eta\), and whether it is the healthy sign depends on the convention for the weight, which~\eqref{eq:Smatrix} does not fix. The Maxwell transformation \(\delta c_i^\nu=\partial^\nu\lambda_i\) is the leading form of \(\delta A^\nu=\{\Lambda,T^\nu\}\) with \(\Lambda=\lambda_ie^i\), under which \(h\) changes by a divergence-free linearised diffeomorphism. Masses, the connection induced by the \(y\)-dependence of \(S_i\) and \(m\), and the mixing of harmonics enter at relative orders \(1/(|q|\ell)\) and \(1/(|q|\ell)^2\), which frozen coefficients do not determine; they are not computed.

This is a statement for the Poisson limit, the stated class with the bilinear term, and the leading symbol. It does not determine masses, the signs of subleading terms or decoupling beyond leading order, and it does not identify the fibre-averaged tensor with the emergent metric of Chapter~\ref{ch:vierbein-emergence}.

\paragraph{Conjugate-coordinate class: what the tree-level equations of the level-one amplitudes imply for \(\mathcal E[h]\) (conditional; bounded no-go).}
Keep the setting of the preceding paragraph. Let \(\mathcal E\) be the linearised Einstein operator of \(\eta\), \(h^{\mu\nu}=q_\rho(b^{\nu;\mu\rho}+b^{\mu;\nu\rho})\) and \(s_i^\mu=S_i^{\mu\rho}q_\rho\), so that \(h^{\mu\nu}=\varkappa\sum_i(s_i^\mu c_i^\nu+s_i^\nu c_i^\mu)\). The order count of the preceding paragraphs becomes precise at the leading symbol. In exact rational arithmetic at the base points \((y^0,y^4)=(5,3),(13,5),(5,4)\) in units \(R=1\) the following hold.

\emph{Free equations.} The equations \(q_a\mathcal F_i^{ab}=0\) make \(c_i\) pure gauge for \(q^2\neq0\), so \(h\) is then a divergence-free linearised diffeomorphism and \(\mathcal E[h]=0\). For null \(q\) in Lorenz gauge there are six physical polarisations (three Maxwell fields): three give \(h=0\), and the other three give \(h\) modulo diffeomorphisms spanned by the two transverse-traceless polarisations and one further polarisation with nonzero trace (base points \((5,3)\) and \((13,5)\)). For these solutions \(\mathcal E[h]_{\mu\nu}=\tfrac12q_\mu q_\nu h^\alpha{}_\alpha\), which vanishes on the transverse-traceless polarisations and not on the polarisation with nonzero trace; the longitudinal diffeomorphism is not a symmetry of the class, so this trace is not removed by gauge.

\emph{No local relation.} The map \(c\mapsto\mathcal E[h(c)]\) is cubic in \(q\), and the Maxwell operator \(c\mapsto q^2c^b-q^b\,q\cdot c\) is quadratic. No matrix \(P(q)\) linear in \(q\) satisfies \(\mathcal E[h(c)]=P(q)\cdot(\text{Maxwell operator on }c)\), and none of the ten components of \(\mathcal E[h(c)]\) is divisible by \(q^2\). Hence no local equation \(\mathcal E[h]=P(\partial)J\) follows from the second-order equations \(\partial_a\mathcal F_i^{ab}=J_i^b\) for a source \(J_i\).

\emph{Response to a source.} Take as input a conserved source, \(q_\mu T^{\mu\nu}=0\), coupled to the amplitudes through \(h\), that is \(J_{i\nu}\propto s_i^\mu T_{\mu\nu}\) (gauge invariance of the coupling requires the conservation). Then \(c_i\) is determined up to gauge, \(c_i^\nu\propto J_i^\nu/q^2\), and
\[
h^{\mu\nu}=\frac{N}{q^2}\bigl(Q^{\mu\alpha}T_\alpha{}^\nu+Q^{\nu\alpha}T_\alpha{}^\mu\bigr),\qquad
Q^{\mu\alpha}=\sum_is_i^\mu s_i^\alpha=q_\perp^2P_\perp^{\mu\alpha}-q_\perp^\mu q_\perp^\alpha ,
\]
with \(N\) a nonzero constant fixed by the couplings, \(P_\perp\) the projector on \(y^\perp\) and \(q_\perp=P_\perp q\), both restricted to the indices \(0,\dots,3\). Modulo \(q\circ\xi\) this is \(h^{\mu\nu}\simeq N\bigl[2(1+q_\parallel^2/q^2)T^{\mu\nu}+\hat y^\mu\tau^\nu+\tau^\mu\hat y^\nu\bigr]\) with \(\hat y=y/\ell\), \(q_\parallel=q_\mu\hat y^\mu\) and \(\tau^\nu=\hat y_\alpha T^{\alpha\nu}\). The response is homogeneous of degree zero in \(q\): \(h\) is algebraic in \(T\), local for \(q_\parallel=0\), in particular for a static source at a base point with \(y^i=0\), and otherwise carries the factor \(q_\parallel^2/q^2\). Since \(\mathcal E\) is of degree two, \(\mathcal E[h]\) is homogeneous of degree two in \(q\) and is nonzero whenever \(h\) is not a linearised diffeomorphism; it therefore cannot equal a constant multiple of \(T\) (the sampled cases also exclude \(q^2T\)). A relation of second order in \(c\) is zeroth order in \(h\) once the common derivative is removed, whereas \(\mathcal E[h]=T\) is second order in \(h\).

\emph{Induced term.} Suppose a term \((16\pi G_{\mathrm{eff}})^{-1}\int\sqrt{-G}\,R\), reported to arise at one loop in~\cite{Steinacker2023} and not computed here, is added as an input. Its quadratic form in \(h\) has degree two in \(q\) and the tree-level form has degree zero, so the tree-level form dominates at small \(q\), and an Einstein-type response \(h\sim G_{\mathrm{eff}}T/q^2\) can appear only above the scale set by the ratio of the two coefficients, provided that scale lies within \(|q|\ell\gg1\); otherwise the induced form is subdominant throughout the regime of the class. Removing the tree-level term would require a loop contribution to the coefficient of \(\sum_i\mathcal F_i^2\) that cancels it; this is not examined.

This is a statement for the Poisson limit, the stated class, the leading symbol, tree level, the level-one sector coupled to a conserved source through \(h\), and the sheet \(y^4>0\). It does not exclude an Einstein-type equation from the induced term, from subleading orders in \(1/(|q|\ell)\), from other sectors or from the entanglement route of Assumption~\ref{ass:EE}, and it derives no value or sign of \(G_{\mathrm{eff}}\). This changes no ledger label; Problem~F remains open.
